% !TEX root = ../main.tex

\chapter{Quasi-smooth derived geometry and Kuranishi charts}
\label[chapter]{chap:Local_Derived_Geometry_Kuranishi_Charts}

The preceding chapters described the functions and tangent complexes of
derived zero loci. We now ask when a derived object admits such a
presentation locally. The answer is detected by its cotangent complex:
the quasi-smooth objects are precisely those locally presented as zero
loci of sections of finite-rank vector bundles.

To make this local statement precise, we first pass from animated smooth
rings to their spectra and structure sheaves. Smooth localization supplies
the restriction maps, and gluing affine objects leads to derived
$\Cinfty$-schemes. We then state the local zero-locus theorem and study
its presentations, called Kuranishi charts. Stabilization and changes of
bundle frame illustrate how different equations can describe the same
derived object. These comparisons will be used later in the finite-dimensional
reduction of elliptic equations.

\section{Smooth localization}
\label[section]{sec:Localization_To_Derived_Spectra}

For a smooth function $a$ on a manifold $M$, put
$D(a)=\{x\in M:a(x)\neq0\}$. The smooth localization satisfies
\[
 \Cinfty(M)[a^{-1}]\cong\Cinfty(D(a)).
\]
One way to see this is to identify $D(a)$ with the graph of $1/a$ in
$M\times\R$. This graph is the transverse zero locus of
$ya(x)-1$, so adjoining a smooth inverse to $a$ gives precisely its
function ring.

Smooth localization is larger than algebraic localization of the underlying
ring. For instance, $\Cinfty(\R)[x^{-1}]$ contains the smooth function
$\exp(1/x^2)$ on the punctured line. It is not a fraction $g(x)/x^n$
with $g$ smooth on all of $\R$: multiplication by any power of $x$
leaves its divergence at zero. Smooth functional calculus is part of
the localization.

Let now $A$ be an animated $\Cinfty$-ring and $a\in\pi_0A$.

\begin{definition}
\label[definition]{def:Derived_Localization}
The smooth localization $A[a^{-1}]$ is initial among animated smooth
rings under $A$ in which the image of $a$ is invertible in $\pi_0$.
\end{definition}

Choose a representative of $a$, hence a map $\Cinfty(\R)\to A$.
The localization is the pushout
\[
 \begin{tikzcd}
 \Cinfty(\R)\rar\dar\drar[pushout]&A\dar\\
 \Cinfty(\R\setminus\{0\})\rar&A[a^{-1}].
 \end{tikzcd}
\]
The universal property makes the result independent of this choice.

\begin{theorem}[Flatness of smooth localization]
\label[theorem]{thm:Homotopy_Groups_Smooth_Localization}
There are natural isomorphisms
\[
 \pi_0(A[a^{-1}])\cong(\pi_0A)[a^{-1}],\qquad
 \pi_i(A[a^{-1}])
 \cong\pi_i(A)\otimes_{\pi_0A}(\pi_0A)[a^{-1}].
\]
The smooth localization of the ordinary ring $\pi_0A$ is flat as
a module over its underlying ring.
\end{theorem}

We use this result without proof
\cite[Proposition~4.1.3.13]{Steffens_Derived_Differential_Geometry}.
The tensor product in its second formula is an ordinary module tensor
product. Flatness explains why restricting functions creates no
additional higher homotopy from a derived tensor product.
It restricts each existing homotopy module.

The corresponding relative cotangent complex vanishes:
\[
 \mathbb L_{A[a^{-1}]/A}\cong0.
\]
By transitivity,
\begin{equation}\label{eq:Cotangent_Complex_Localizes}
 \mathbb L_{A[a^{-1}]}\cong\mathbb L_A\otimes_A A[a^{-1}].
\end{equation}
This is the compatibility needed to regard cotangent complexes as local
geometric objects.

\section{The spectrum and its structure sheaf}

Assume $A$ is finitely presented as an animated smooth ring.
Its spectrum has underlying set
\[
 |\Spec A|=\Hom_{\Cinfty\text{-rings}}(\pi_0A,\R).
\]
These are real points, not arbitrary prime ideals of an underlying
commutative ring. A point determines a maximal ideal with residue field
$\R$, namely the kernel of its evaluation map. The evaluation map
is the useful datum in smooth geometry.

For $a\in\pi_0A$, the subsets
\[
 D(a)=\{x:x(a)\neq0\}
\]
form a basis for a topology, since $D(a)\cap D(b)=D(ab)$.
If $\pi_0A=\Cinfty(\R^n)/I$, evaluation identifies this topological
space with the common zero set $Z(I)\subseteq\R^n$, with its subspace
topology. To check the topology, bump functions give basic
neighbourhoods inside any prescribed Euclidean open set.

\begin{construction}[Derived spectrum]
\label[construction]{con:Derived_Spectrum}
The derived spectrum is
\[
 \Spec A=(|\Spec A|,\mathcal O_{\Spec A}),
\]
where the sheaf of animated smooth rings is obtained from smooth
localizations. In the finite-presentation setting,
\[
 \mathcal O_{\Spec A}(D(a))\cong A[a^{-1}].
\]
The local ring at $x$ is the filtered colimit of these localizations
over basic neighbourhoods of $x$.
\end{construction}

The assertion that these assignments have the required sheaf and
reconstruction properties is a theorem about smooth rings.
It is part of the affine comparison
\cite[Section~4.1.3, especially Corollary~4.1.3.34]
{Steffens_Derived_Differential_Geometry}.
For unrestricted smooth rings one must take care with sheafification
and reconstruction; finite presentation is the class used here.

The term \emph{locally animated $\Cinfty$-ringed space} means that
$\pi_0$ of every stalk is a local smooth ring. At $x$ its residue map
is evaluation in $\R$. The higher homotopy modules remain in the stalk.
A stalk belongs to a point; the values on basic opens are the
localizations before taking that colimit.

For an ordinary manifold, the construction recovers its familiar
ringed space:
\[
 \Spec\Cinfty(M)\cong(M,\Cinfty_M).
\]
Real homomorphisms are evaluations at points, the basic opens give
the manifold topology, and the localization formula identifies the
structure sheaf with smooth functions.

The examples from \Cref{chap:Derived_Zero_Loci_And_Tangent_Complexes} now have an explicit topological meaning.
The spectra of $\R$, $\R[\delta]/(\delta^2)$, and
$\R[\varepsilon]$ with $|\varepsilon|=1$, $d=0$, all have one point.
Their structure sheaves differ. The second has nilpotents in degree zero;
the third has a nonzero homotopy group in degree one.

\section{Derived schemes and quasi-smoothness}

The spectrum construction supplies affine objects and their open subsets.
We can therefore define locally affine objects, then use the cotangent
complex to single out those that should admit local presentations by
vector bundle sections.

An open subset $V\subseteq|X|$ of a locally ringed object defines
an open subobject $(V,\mathcal O_X|_V)$. In particular,
\[
 \Spec A[a^{-1}]
 \cong(D(a),\mathcal O_{\Spec A}|_{D(a)}).
\]
Thus algebraic localization becomes geometric open restriction.

\begin{definition}
\label[definition]{def:Derived_Cinfty_Scheme}
A derived $\Cinfty$-scheme locally of finite presentation is a locally
animated $\Cinfty$-ringed space with an open cover by spectra of
finitely presented animated $\Cinfty$-rings.
\end{definition}

The affine objects in this definition are exactly the derived manifolds
of the universal-property convention in \Cref{chap:From_Smooth_Algebra_To_Derived_Manifolds}:
\[
 \mathrm{DMfd}\simeq
 \bigl(\operatorname{Alg}_{\Cinfty}(\An)^{\mathrm{fp}}\bigr)^{\mathrm{op}}.
\]
In particular, an ordinary manifold is affine in this sense.
Affineness does not mean that it is a single Euclidean coordinate chart.
For example, $S^1=\Spec\Cinfty(S^1)$ is affine.
Allowing locally affine objects supplies the category in which
open gluing can be performed.

The classical truncation is
\[
 t_0X=(|X|,\pi_0\mathcal O_X).
\]
It keeps the topology and removes higher homotopy from the structure
sheaf. It can still contain nilpotents. Thus classical truncation and
the underlying zero set are different operations.

Cotangent complexes restrict by
\Cref{eq:Cotangent_Complex_Localizes}, so they glue over affine opens.
A module over an animated ring $A$ is \emph{perfect} if it is a retract
of a finite construction by cofibers from shifts of finite sums of $A$.
For an ordinary ring this is equivalent to having a bounded complex
of finitely generated projective modules as a presentation. Perfection
on a derived scheme is checked on its affine opens. A module has Tor-amplitude
in $[a,b]$ if tensoring with any discrete module produces homology
only in degrees $a,\ldots,b$.
This condition concerns derived tensor products; it is stronger than
merely locating the homology of the original complex.

\begin{definition}[Quasi-smoothness]
\label[definition]{def:Quasi_Smooth_Derived_Cinfty_Scheme}
A derived $\Cinfty$-scheme locally of finite presentation is
\emph{quasi-smooth} if its cotangent complex is perfect of
Tor-amplitude contained in $[0,1]$.
For a morphism, the same definition uses the relative cotangent complex.
\end{definition}

Perfection and Tor-amplitude can be checked locally.
A two-term complex of finite-rank vector bundles in degrees $1,0$
has this amplitude. The cotangent calculation of \Cref{chap:Stable_Linearization_And_Tangent_Complexes} therefore
proves that every derived zero locus of a section of a finite-rank
bundle is quasi-smooth.

For a perfect complex the amplitude condition can also be tested on
residue fields. Consequently, at each real point of a quasi-smooth
object its tangent complex has homology only in degrees $0,-1$.
The bundle presentation of the cotangent complex, rather than just
a dimension count at individual points, is the local structure behind
this assertion.

\section{Local presentations by Kuranishi charts}
\label[section]{sec:Kuranishi_Charts_As_Presentations}

The cotangent calculation already shows that a vector bundle zero locus
is quasi-smooth. The local normal-form theorem supplies the converse.
We first specify what data is included in a chart.

\begin{definition}[Kuranishi chart]
\label[definition]{def:Kuranishi_Chart}
A Kuranishi presentation is a triple $(U,E,s)$ consisting of a
finite-dimensional manifold, a finite-rank vector bundle on it,
and a smooth section. Its associated derived object is
$Z^{\mathrm{der}}(s)$.
A Kuranishi chart for $X$ also includes an identification of this
derived zero locus with an open subobject of $X$.
\end{definition}

This usage concerns charts without isotropy, boundary, or corners.
It is the local presentation needed in the representability proof.
A choice of charts alone does not specify any of the more elaborate
global notions called Kuranishi structures in the literature.

\begin{theorem}[Local zero-locus form]
\label[theorem]{thm:Local_Zero_Locus_Form}
A derived $\Cinfty$-scheme locally of finite presentation is
quasi-smooth if and only if it admits an open cover by Kuranishi charts.
\end{theorem}

The forward implication is the substantive normal-form theorem.
The affine statement is proved in
\citet[Proposition~5.1.1.11 and Corollary~5.1.1.13]
{Steffens_Derived_Differential_Geometry}.
In that statement, a finitely presented quasi-smooth affine object
can be expressed as the zero locus of a bundle section over an open
subset of a Euclidean space. Applying it on affine opens proves the
displayed local form. The reverse implication is the cotangent calculation.

The proof begins by choosing generators for the ring, giving an
ambient smooth affine object. The relative cotangent complex describes
its equations. The amplitude assumption implies that the shifted
relative cotangent complex is a finite projective module, which can
be extended to a vector bundle on an ambient neighbourhood.
Lifting the relations then produces a section of that bundle.
A cotangent criterion identifies its derived zero locus with the
original object. This explains the role of the amplitude assumption;
the extension and comparison theorems in this argument are quoted.

At a point of a chart, the tangent complex is
$[T_xU\to E_x]$, and the virtual dimension is
$\dim U-\operatorname{rank}E$. Quasi-smoothness permits both singular
classical truncations and nonzero higher homotopy. The zero locus of
$x^2$ illustrates the first, and a zero equation illustrates the second.

\section{Changing a presentation}

The quotient rings, Koszul algebras, and tangent complexes calculated
earlier distinguish different geometric objects. We now use those
calculations to recognize changes of equations that preserve an object.

Let $f\colon U\to\R^r$, and define
\[
 f'\colon U\times\R^q\longrightarrow\R^r\times\R^q,
 \qquad f'(x,y)=(f(x),y).
\]
The additional equation forces the additional variable to zero.

\begin{proposition}[Stabilization]
\label[proposition]{prop:Stabilization_Kuranishi_Chart}
Projection onto $U$ induces an isomorphism
$Z^{\mathrm{der}}(f')\cong Z^{\mathrm{der}}(f)$.
\end{proposition}

\begin{proof}
The second factor is the derived zero locus of the identity of $\R^q$.
That map is transverse to zero, so its zero locus is the ordinary point.
Pasting pullbacks gives the stated isomorphism.
Equivalently, the new Koszul variables satisfy $d\eta_j=y_j$.
Evaluation at $y=\eta=0$ removes an acyclic Koszul resolution.
\end{proof}

For example, $x^2$ and $(x^2,y)$ give the same derived object.
Their tangent presentations differ by
$[\R\xrightarrow1\R]$, an acyclic summand.
Both ambient dimension and the number of equations increase by one,
leaving virtual dimension unchanged.

Changing a bundle frame likewise changes the list of equations
without changing the object. If $g(x)$ is an invertible matrix of
smooth functions, then $f$ and $g f$ present isomorphic derived zero loci.
Explicitly, the map from the Koszul algebra for $gf$ to that for $f$
sends its generator $\varepsilon'_i$ to $\sum_jg_{ij}\varepsilon_j$.
The inverse matrix gives the inverse dg-algebra map.
The intrinsic bundle model with differential $\iota_s$ makes this
independence immediate.

In contrast, $x$ and $(x,0)$ have the same ordinary zero set but
different derived zero loci. The extra zero equation contributes a
degree-minus-one tangent direction and a degree-one homotopy class of functions.
Nor does equality of tangent complexes suffice: $x^2$ and $x^3$
have the same tangent complex at zero but distinct classical
function rings. Comparisons of charts must respect their equations.

\section{Further reading}

Two refinements of the chart description will recur later: removing
transverse directions makes a chart minimal at a point, while a
nontrivial obstruction bundle shows what local frames can conceal.

A useful refinement makes a chart minimal at a chosen point.
Suppose $f\colon U\subseteq\R^n\to\R^r$ vanishes at $x$ and $d_xf$
has rank $q$. After linear changes of coordinates, $q$ components
of $f$ have independent differentials. The inverse function theorem
allows these components to be used as coordinates $y\in\R^q$.
Writing the remaining coordinates as $z$, the map becomes
\[
 (y,z)\longmapsto(y,g(y,z)).
\]
Hadamard's lemma gives a smooth matrix $H$ with
$g(y,z)=g(0,z)+H(y,z)y$.
An invertible change of equation generators subtracts $H(y,z)y$,
giving $(y,g(0,z))$. Stabilization now removes the coordinates $y$.

The resulting local chart is the zero locus of
\[
 \kappa\colon\R^{n-q}\supseteq V\longrightarrow\R^{r-q},
 \qquad\kappa(z)=g(0,z),
\]
with $d\kappa$ zero at the chosen point.
Its source tangent and target identify with
$\ker d_xf$ and $\operatorname{coker}d_xf$.
These identifications use choices; the derived object they present is
the original open neighbourhood. The analytic reduction of an elliptic
equation will produce a chart of this kind from an initially
infinite-dimensional problem.

\begin{example}[An obstruction bundle with no global frame]
Let $L\to S^1$ be the M\"obius line bundle and form
$X=S^1\times_L S^1$ using two zero sections.
The intrinsic Koszul algebra is the exterior algebra of $L^\vee$
with zero differential. Therefore
\[
 \pi_0\mathcal O_X=\Cinfty_{S^1},\qquad
 \pi_1\mathcal O_X=\Gamma(-,L^\vee).
\]
Each trivializing open gives a single zero equation. On overlaps the
generator changes by the transition function of $L^\vee$.
There is no globally nonvanishing generator, although the derived
self-intersection is defined globally.

After restriction to the classical circle, the tangent complex is
$[TS^1\xrightarrow0 L]$. Its obstruction sheaf remembers the
nontrivial line bundle. This shows why vector bundles, rather than only
fixed vector spaces of equations, naturally occur in Kuranishi charts.
\end{example}

\begin{exercise}
Show that the derived zero locus of $(x,xy)\colon\R^2\to\R^2$
is locally, and in this case globally, the product of a line and
the derived zero locus of the zero map from a point to $\R$.
\end{exercise}

\begin{exercise}
For $f(x,y)=x^2+y^2$, calculate the tangent complex at the origin
and compare its virtual dimension with the dimension of its zero set.
Explain why the zero set alone cannot carry this information.
\end{exercise}

For ordinary smooth spectra and their structure sheaves, see
\cite[Sections~2.2 and~3.1]{Joyce_Introduction_Cinfty_Schemes}.
The localization and normal-form statements used here are in
\cite[Sections~4.1.3 and~5.1.1]{Steffens_Derived_Differential_Geometry}.
They provide the passage from the affine universal theory to the
local charts used in geometric moduli problems.
